\section{任意正模数的矩阵树计数}
实现见第~\pageref{compact-MatrixTreeMod}~页。
拉普拉斯余子式表示生成树权值和的恒等式本身不要求模数为素数；
需要换掉的是依赖模逆元的行列式算法。
本节沿用无向、内向与外向三种构造，但使用欧几里得行列式消元。
输入边权和正模数均支持 signed 64 位，接口名为
\texttt{MatrixTreeMod::count}。

先把边权正规化到 $[0,m)$，每次矩阵元素加减仍先扩展至 128 位，
再取模存回 64 位。重边的权值和可能超过 64 位范围，不能先累加到普通整数中。
只构造删除根后的余子式，并将其存储直接交给行列式函数。
单点返回 $1\bmod m$，模 1 返回零。
负权与零权仍按代数权值处理，返回零不能直接解释为没有生成树。

独立参考程序逐个枚举边集，用任意精度整数计算权值和。
已检查多个合数模数、所有根与三种方向、断图、重边自环及 64 位极值权重；
另以大权值完全图的 Cayley 公式验证 160 点规模。
普通与 ASan/UBSan 均通过，尚无该接口的在线 AC。
此前 P6178 的 AC 属于素数模版本，不自动转移到本接口。

